% Source: Zhen Liu and Qinghou Zeng, arXiv:2610.15539v1, Section 1.
% Definitions, Tuza conjecture and baseline edge-count theorem below are verbatim source excerpts.
\documentclass{article}
\usepackage{amssymb,amsmath,amsthm}
\newtheorem{theorem}{Theorem}[section]
\newtheorem{conjecture}[theorem]{Conjecture}
\begin{document}
All graphs considered in this paper are finite, simple, and undirected. A \emph{triangle} is a cycle of length three. A \emph{triangle packing} is a family of pairwise edge-disjoint triangles, while a \emph{triangle cover} is a set of edges meeting every triangle. We denote by $\nu(G)$ the maximum size of a triangle packing in $G$, and by $\tau(G)$ the minimum size of a triangle cover. The union of the edges of a maximum triangle packing is a triangle cover, and hence
$
\nu(G)\le \tau(G)\le 3\nu(G).
$
\begin{conjecture}[Tuza]\label{conj:tuza}
Every graph $G$ satisfies
$
\tau(G)\le 2\nu(G).
$
\end{conjecture}
\begin{theorem}\label{thm:main}
For every fixed real number $c<2$, every sufficiently large $n$-vertex graph $G$ with
$
e(G)\ge \frac{n^2}{4}-cn
$
satisfies
$
\tau(G)\le2\nu(G).
$
\end{theorem}
\end{document}
